\documentclass[11pt]{article}
\usepackage[utf8]{inputenc}
\usepackage[margin=1in]{geometry}
\usepackage{amsmath}
\usepackage{graphicx}
\usepackage{booktabs}
\usepackage{hyperref}
\usepackage{natbib}
\usepackage{lineno}
\usepackage{xcolor}
\usepackage{multirow}
\usepackage{float}

\linenumbers

\title{Dose-Dependent Neural Responses in a Thalamic Circuit Driven by Beliefs About Nicotine Strength in Human Smokers}

\author{
Author A$^{1,*}$, Author B$^{1}$, Author C$^{2}$, Author D$^{1,2}$ \\
\\
$^{1}$Department of Neuroscience, Virginia Tech \\
$^{2}$Fralin Biomedical Research Institute \\
$^{*}$Corresponding author: authora@vtc.vt.edu
}

\date{}

\begin{document}
\maketitle

\begin{abstract}
Beliefs profoundly influence behavior and wellbeing, yet the neural mechanisms by which subjective beliefs modulate brain circuits---particularly in a graded, dose-dependent manner---remain poorly understood. Here, we tested whether beliefs about the nicotine content of an e-cigarette can produce dose-dependent neural responses in the thalamus, a region with among the highest densities of nicotinic acetylcholine receptors (nAChRs) in the human brain. Twenty nicotine-dependent adults completed three fMRI sessions in which they were instructed that an e-cigarette contained ``low,'' ``medium,'' or ``high'' nicotine, while actual nicotine content was held constant at 1.2\%. The belief manipulation successfully modulated perceived nicotine strength in a graded manner (3.52, 4.52, 5.82 on a 10-point scale; repeated-measures ANOVA $F_{2,38} = 9.71$, $p = 0.0004$, $\eta_p^2 = 0.34$), while actual nicotine intake ($F_{2,38} = 1.81$, $p = 0.178$), cotinine metabolism ($F_{2,32} = 0.96$, $p = 0.393$), and baseline CO levels ($F_{2,32} = 0.36$, $p = 0.698$) did not differ across conditions. Thalamic activity during a value-based decision-making task increased parametrically with belief dosage (cluster-level $p_{\mathrm{FWE}} = 0.006$; ROI analysis $p = 0.036$), while striatal (nucleus accumbens) reward signals showed no belief modulation ($p = 0.945$). Psychophysiological interaction analysis revealed that thalamus--ventromedial prefrontal cortex (vmPFC) functional connectivity also scaled parametrically with belief dosage. These findings demonstrate that subjective beliefs can produce graded, dose-dependent neural responses in a circuit known to be pharmacologically sensitive to nicotine, establishing a neural dose-response function for beliefs.
\end{abstract}

\section{Introduction}

The influence of beliefs on neural processing is well documented in the placebo literature, where expectations about treatment can modulate pain \citep{wager2004, petrovic2002, atlas2012}, drug responses \citep{volkow2003}, and reward processing \citep{gu2015}. However, most prior studies have examined beliefs as binary variables---active versus placebo---leaving open the question of whether beliefs can produce graded, dose-dependent neural modulation analogous to pharmacological dose-response relationships.

The thalamus is an ideal candidate region for investigating belief-driven dose-response effects in the context of nicotine. It contains among the highest densities of nAChRs in the human brain, particularly $\alpha_4\beta_2$ receptors \citep{brody2006, picciotto2012}. Positron emission tomography (PET) studies have shown that habitual smoking can produce substantial nAChR occupancy in the thalamus even with minimal nicotine exposure \citep{brody2006}, and thalamic responses to actual nicotine follow dose-dependent curves. If beliefs about nicotine strength can modulate thalamic circuits in a similarly graded fashion, this would suggest a deep integration between cognitive representations and pharmacologically sensitive neural substrates.

Here, we addressed this question using a within-subject design in which nicotine-dependent smokers completed three fMRI sessions with identical nicotine exposure but different instructed beliefs about nicotine content (``low,'' ``medium,'' ``high''). We measured reward-related neural activity during a value-based decision-making task \citep{bartra2013} and asked whether belief dosage parametrically modulated thalamic responses and thalamus--vmPFC connectivity. We also tested the anatomical specificity of any belief effect by examining the nucleus accumbens, a key reward processing region \citep{haber2006}.

\section{Results}

\subsection{Participant Characteristics}

Twenty nicotine-dependent adults (14 male, 6 female; mean age 41.1 $\pm$ 12.0 years; range 24--61; $\geq$10 cigarettes/day for $>$1 year) completed all three sessions (Table~\ref{tab:participants}). An independent healthy control cohort (31 participants; 16 male, 15 female; mean age 28 $\pm$ 9 years) performed the same task without nicotine manipulation.

\begin{table}[H]
\centering
\caption{Participant demographics and inclusion criteria. Three participants were excluded from the original 23 enrolled (1 software malfunction, 1 fell asleep during scanning, 1 lost behavioral data). The healthy control cohort had 2 exclusions for excessive head movement ($>$3~mm). Power analysis based on prior work \citep{gu2015} yielded a required $N = 18$--$20$ (Cohen's $d = 0.69$; 90\% power at $\alpha = 0.05$).}
\label{tab:participants}
\begin{tabular}{lcc}
\toprule
\textbf{Characteristic} & \textbf{Smokers ($n=20$)} & \textbf{Controls ($n=31$)} \\
\midrule
Sex (F/M) & 6/14 & 15/16 \\
Age (mean $\pm$ SD) & 41.1 $\pm$ 12.0 & 28 $\pm$ 9 \\
Age range (years) & 24--61 & 19--52 \\
Handedness & All right & All right \\
Cigarettes/day ($\geq$) & 10 & 0 (exclusion) \\
Prior e-cigarette use & None & N/A \\
Sessions per participant & 3 & 1 \\
Session spacing & $\sim$1 week & N/A \\
\bottomrule
\end{tabular}
\end{table}

\subsection{Belief Manipulation Succeeded While Pharmacological Exposure Was Constant}

The instructed belief manipulation produced a graded increase in perceived nicotine strength while all pharmacological and physiological control measures remained constant across conditions (Table~\ref{tab:sanity_checks}).

\begin{table}[H]
\centering
\caption{Manipulation checks and control measures across belief conditions. Values are mean $\pm$ SD. Repeated-measures ANOVA $F$-statistics, $p$-values, and effect sizes ($\eta_p^2$) are reported. Only perceived nicotine strength differed significantly, confirming successful belief manipulation with constant pharmacological exposure. 90\% CIs are reported for $\eta_p^2$.}
\label{tab:sanity_checks}
\begin{tabular}{lcccccc}
\toprule
\textbf{Measure} & \textbf{Low} & \textbf{Medium} & \textbf{High} & \textbf{$F$} & \textbf{$p$} & \textbf{$\eta_p^2$ [90\% CI]} \\
\midrule
Perceived strength (AU) & 3.52 $\pm$ 0.61 & 4.52 $\pm$ 0.41 & 5.82 $\pm$ 0.47 & 9.71 & \textbf{0.0004} & 0.34 [0.12, 0.48] \\
Nicotine intake (mg) & 0.93 $\pm$ 0.56 & 0.72 $\pm$ 0.42 & 0.78 $\pm$ 0.43 & 1.81 & 0.178 & 0.09 [0.00, 0.22] \\
$\Delta$Cotinine (ng/mL) & 14.3 $\pm$ 8.7 & 12.8 $\pm$ 7.2 & 13.6 $\pm$ 8.1 & 0.96 & 0.393 & 0.06 [0.00, 0.17] \\
Exhaled CO (ppm) & 18.4 $\pm$ 6.3 & 17.9 $\pm$ 5.8 & 18.1 $\pm$ 6.1 & 0.36 & 0.698 & 0.02 [0.00, 0.10] \\
\bottomrule
\end{tabular}
\end{table}

\subsection{Thalamic Activity Increases Parametrically with Belief Dosage}

Whole-brain analysis revealed a significant cluster in the thalamus where reward-tracking signals (parametric modulation by market return) differed across belief conditions (Table~\ref{tab:thalamus}).

\begin{table}[H]
\centering
\caption{Thalamic activation analysis. Top: whole-brain repeated-measures ANOVA identifying belief-modulated clusters ($k \geq 50$ voxels, cluster-level FWE correction). Bottom: ROI analysis using an independent anatomical thalamus mask. Parameter estimates reflect reward-tracking signal strength. Healthy control data shown for comparison. Trend test: linear contrast across Low $<$ Medium $<$ High.}
\label{tab:thalamus}
\begin{tabular}{lccccc}
\toprule
\textbf{Analysis} & \textbf{Low} & \textbf{Medium} & \textbf{High} & \textbf{Statistic} & \textbf{$p$-value} \\
\midrule
\multicolumn{6}{l}{\textit{Whole-brain cluster (thalamus)}} \\
Cluster size (voxels) & \multicolumn{3}{c}{---} & $k = 127$ & $p_{\mathrm{FWE}} = 0.006$ \\
\midrule
\multicolumn{6}{l}{\textit{ROI analysis (independent mask)}} \\
$\beta$ estimate (AU) & 0.42 $\pm$ 0.18 & 0.67 $\pm$ 0.21 & 0.89 $\pm$ 0.24 & $F_{2,38} = 3.62$ & 0.036 \\
Linear trend & \multicolumn{3}{c}{Low $<$ Medium $<$ High} & $t_{19} = 2.74$ & 0.013 \\
Controls (AU) & \multicolumn{3}{c}{0.71 $\pm$ 0.19} & --- & --- \\
\bottomrule
\end{tabular}
\end{table}

Thalamic reward-tracking activity increased parametrically from Low to Medium to High belief conditions, with a significant linear trend ($t_{19} = 2.74$, $p = 0.013$). Notably, the High belief condition produced thalamic activity exceeding that of healthy controls, while the Low condition produced activity below the control level.

\subsection{Striatum Shows No Belief Modulation (Anatomical Specificity)}

To establish anatomical specificity, we examined the nucleus accumbens (NAcc), a key reward processing structure (Table~\ref{tab:striatum}).

\begin{table}[H]
\centering
\caption{Nucleus accumbens ROI analysis. The NAcc showed robust reward-tracking signals collapsed across conditions but no parametric modulation by belief dosage. Permutation test ($n = 5{,}000$) confirmed the null result. This demonstrates anatomical specificity of the belief dose-response to the thalamus.}
\label{tab:striatum}
\begin{tabular}{lccccc}
\toprule
\textbf{Region} & \textbf{Low} & \textbf{Medium} & \textbf{High} & \textbf{$F_{2,38}$} & \textbf{$p$-value} \\
\midrule
NAcc $\beta$ (AU) & 1.12 $\pm$ 0.34 & 1.14 $\pm$ 0.31 & 1.11 $\pm$ 0.36 & 0.06 & 0.945 \\
NAcc (permutation) & --- & --- & --- & --- & 0.940 \\
Controls NAcc (AU) & \multicolumn{3}{c}{1.08 $\pm$ 0.28} & --- & --- \\
\midrule
Collapsed (all conditions) & \multicolumn{3}{c}{1.12 $\pm$ 0.33} & $t_{19} = 4.82$ & $<$0.001 \\
\bottomrule
\end{tabular}
\end{table}

The NAcc showed robust reward-tracking signals ($t_{19} = 4.82$, $p < 0.001$ collapsed across conditions) but no modulation by belief dosage ($F_{2,38} = 0.06$, $p = 0.945$), confirmed by non-parametric permutation testing ($p = 0.940$). This double dissociation---thalamic but not striatal belief modulation---establishes anatomical specificity.

\subsection{Thalamus--vmPFC Connectivity Scales with Belief Dosage}

Psychophysiological interaction (PPI) analysis \citep{friston1997} with the thalamus as seed region revealed belief-dependent changes in functional connectivity with vmPFC (Table~\ref{tab:ppi}).

\begin{table}[H]
\centering
\caption{Psychophysiological interaction analysis: thalamus--vmPFC functional connectivity as a function of belief dosage. PPI $\beta$ values reflect the interaction between thalamic activity and task context modulated by belief condition. Linear trend test across conditions. Cluster-based permutation test for the PPI contrast ($n = 5{,}000$ permutations).}
\label{tab:ppi}
\begin{tabular}{lcccc}
\toprule
\textbf{Connectivity} & \textbf{Low} & \textbf{Medium} & \textbf{High} & \textbf{$p$-value} \\
\midrule
Thalamus--vmPFC PPI $\beta$ (AU) & 0.08 $\pm$ 0.12 & 0.19 $\pm$ 0.14 & 0.34 $\pm$ 0.16 & 0.008 \\
Linear trend & \multicolumn{3}{c}{Low $<$ Medium $<$ High} & 0.004 \\
Permutation test & \multicolumn{3}{c}{---} & 0.011 \\
\bottomrule
\end{tabular}
\end{table}

Thalamus--vmPFC coupling increased parametrically with belief dosage ($p = 0.008$; linear trend $p = 0.004$), revealing a circuit-level dose-response relationship. The vmPFC has been implicated in value representation and state inference \citep{bartra2013}, suggesting that belief-dependent thalamic modulation is communicated to prefrontal circuits involved in subjective valuation.

\subsection{Subjective Perception Correlates with Thalamic Activity}

Across individuals, self-reported perceived nicotine strength correlated with thalamic reward-tracking activity (Pearson $r = 0.48$, $p = 0.032$, $n = 20$), bridging the gap between subjective experience and neural representation. This correlation was specific to the thalamus; no significant correlation was found for the NAcc ($r = 0.04$, $p = 0.872$).

\section{Discussion}

Our findings demonstrate that subjective beliefs about nicotine strength can produce graded, dose-dependent neural responses in the thalamus---a region with exceptionally high nAChR density---despite constant actual nicotine exposure. This extends prior binary (active/placebo) belief studies \citep{gu2015, volkow2003} to a dose-response framework, suggesting that the brain constructs parametric neural representations of belief strength.

The thalamic locus of the belief effect is particularly noteworthy given the region's pharmacological profile. The thalamus contains the highest density of $\alpha_4\beta_2$ nAChRs in the human brain \citep{brody2006}, and PET studies have shown that even minimal nicotine exposure can produce substantial thalamic nAChR occupancy in habitual smokers. Our results suggest that beliefs about nicotine strength may modulate the same thalamic circuits that respond to actual nicotine, potentially through top-down predictions about expected pharmacological effects \citep{friston2010}.

The anatomical specificity of the belief effect---present in the thalamus but absent in the nucleus accumbens---constrains mechanistic accounts. The NAcc showed robust reward-tracking signals but was entirely insensitive to belief dosage, suggesting that beliefs do not globally modulate the reward system but rather selectively influence circuits rich in the molecular targets (nAChRs) relevant to the belief content \citep{picciotto2012, dani2011}. This specificity also argues against a trivial explanation based on general arousal or motivational differences between belief conditions.

The parametric scaling of thalamus--vmPFC connectivity with belief dosage reveals a circuit-level mechanism for belief-driven neural modulation. The vmPFC has been consistently implicated in value representation \citep{bartra2013} and state inference, suggesting that belief-dependent thalamic modulation is integrated into prefrontal computations of subjective value.

Several limitations should be noted. First, the sample size ($n = 20$), while adequately powered based on prior work \citep{gu2015, faul2007}, limits the detection of smaller effects. Second, the absence of a nicotine-free placebo condition means we cannot fully disentangle belief effects from pharmacological effects. Third, the healthy control cohort differed in age from the smoker cohort, precluding direct statistical comparison. Future work should include age-matched controls and a fully balanced design crossing belief with actual nicotine dosage.

\section{Methods}

\subsection{Participants}

Twenty-three nicotine-dependent adults were enrolled; three were excluded (1 software malfunction, 1 fell asleep, 1 lost behavioral data), yielding a final sample of 20. Inclusion criteria: $\geq$10 cigarettes/day for $>$1 year, right-handed, no prior e-cigarette use, not currently attempting to quit, no illicit drug use in the past 2 months. A separate healthy control cohort of 33 participants (2 excluded for $>$3~mm head movement) performed the same task.

\subsection{Experimental Procedure}

Each participant completed three fMRI sessions spaced approximately one week apart. In each session: (1) saliva sample collected (pre-vaping cotinine); (2) exhaled CO measured; (3) belief condition instructed (``low,'' ``medium,'' or ``high'' nicotine); (4) 20 minutes of controlled vaping with a commercially available e-cigarette (1.2\% nicotine, constant across all conditions); (5) saliva sample collected (post-vaping); (6) fMRI scan during value-based decision-making task. Condition order was counterbalanced across participants.

\subsection{Biochemical Measures}

Nicotine intake was computed as the change in e-cigarette cartridge weight multiplied by 1.2\% nicotine content. Cotinine was quantified from saliva using high-performance liquid chromatography tandem mass spectrometry (LC-MS/MS). Exhaled CO was measured with a calibrated CO monitor.

\subsection{fMRI Acquisition and Analysis}

Functional images were acquired at 3T (TR = 2~s, voxel size = 3~mm isotropic). Preprocessing followed standard SPM12 pipelines \citep{woolrich2009}: realignment, slice-timing correction, normalization to MNI space, and 6~mm FWHM smoothing. First-level models included parametric modulation by market return (value-tracking regressor) and motion parameters. Second-level repeated-measures ANOVA tested for belief effects. Cluster-level FWE correction ($k \geq 50$ voxels) was applied \citep{winkler2014}. ROI analyses used independent anatomical masks for thalamus and nucleus accumbens.

\subsection{PPI Analysis}

PPI analysis \citep{friston1997} used the thalamic cluster as seed region. The interaction between seed time course and task contrast (parametric modulation by market return) was computed for each belief condition, and the belief effect on PPI coefficients was tested with repeated-measures ANOVA.

\subsection{Statistical Analysis}

Repeated-measures ANOVA was used for within-subject comparisons across three belief conditions, with Greenhouse--Geisser correction when sphericity was violated. Effect sizes are reported as partial eta-squared ($\eta_p^2$) with 90\% confidence intervals. Trend tests used linear contrasts. Non-parametric permutation tests ($n = 5{,}000$) were conducted to confirm parametric results. All tests were two-tailed with $\alpha = 0.05$.

\section{Conclusion}

This study establishes that subjective beliefs about nicotine strength can produce dose-dependent neural responses in the thalamus, a region rich in nicotinic acetylcholine receptors, despite constant actual nicotine exposure. The anatomical specificity (thalamus but not striatum) and circuit-level scaling (thalamus--vmPFC connectivity) suggest that beliefs engage the same molecular and circuit substrates targeted by actual pharmacological agents. These findings extend the placebo literature from binary to parametric belief effects and have implications for understanding how expectations shape drug responses, addiction, and potentially therapeutic interventions.

\bibliographystyle{apalike}
\bibliography{references}

\end{document}
